\documentclass[11pt,a4paper]{article}

\usepackage[margin=2.3cm]{geometry}
\usepackage[T1]{fontenc}
\usepackage{lmodern}
\usepackage{amsmath}
\usepackage{booktabs}
\usepackage{tabularx}
\usepackage{longtable}
\usepackage{enumitem}
\usepackage{xcolor}
\usepackage{listings}
\usepackage{hyperref}

\hypersetup{colorlinks=true, linkcolor=blue!50!black, urlcolor=blue!50!black}
\lstset{basicstyle=\ttfamily\small, columns=fullflexible, breaklines=true, frame=single,
  framerule=0.3pt, rulecolor=\color{black!30}, backgroundcolor=\color{black!3}, showstringspaces=false}
\newcommand{\code}[1]{\texttt{#1}}

\title{CPU--GPU Co-Scheduling Module for CloudSim 7G\\[4pt]
\large Phase 4 Report: Validation against Synergy's Simulator}
\author{Chirag \and Hamiz}
\date{23 November 2026}

\begin{document}
\maketitle

\begin{abstract}
Phase~4 is the gate for every later result: our module must reproduce Synergy's published simulator. We ran
Synergy's simulator ourselves for 31 configurations (the reference run, a load sweep, four orderings with
Greedy, Tune and the fixed CPU share, three CPU:GPU ratios, three workload splits and multi-GPU jobs), recorded
every job each run generated, replayed those jobs in our simulator, and compared every measured job. In all
31 configurations every one of the 1{,}000 measured jobs gets exactly the job completion time Synergy gives it,
to the second; the reference run's mean of 21.8766\,h and the fixed share's 81.0886\,h (the paper's
``81\,hrs'') are reproduced to every printed digit. The pass rule set in advance was a mean within 10\%.
Getting there found three real bugs in our module, all fixed and covered by tests, and two issues outside it.
\end{abstract}

\tableofcontents

%======================================================================
\section{Method}
%======================================================================

\subsection{Running Synergy's simulator}
Synergy's simulator (commit \code{1bb2bed}) was run on this machine (Python~3.14) on private copies of its code,
so the repository checkout is untouched. Four compatibility changes were needed, none of which touch scheduling:
\code{collections.Mapping} to \code{collections.abc.Mapping}, matplotlib's \code{basex} to \code{base}, gRPC stubs
generated from the \code{.proto} files exactly as Synergy's \code{deployment/Makefile} does, and a stand-in
\code{numa} module (NUMA is used only for physical deployment; the stand-in raises if ever called). Synergy also
needs a \code{cache/} folder to save its results. The reference run reproduces the Check~2 reference file from
another machine exactly (1000/1000 jobs, mean 21.8766\,h), so the environment does not change results.

\subsection{Replaying Synergy's own jobs}
Synergy can record every job it generates (\code{--record\_trace}): id, model, arrival, iterations and GPUs,
including jobs still unfinished when it stops after the measured jobs (ids 4000--4999) are done. A new tool,
\code{tools/synergy\_record\_to\_jobs.py}, turns that record into our job CSV; Synergy's exact per-job JCTs are read
from its result file. Each configuration's replay uses the trace recorded by \emph{that same run}, and its
experiment file sets every option from Synergy's code: 300\,s rounds with Synergy's whole-iteration progress
(\code{synergy\_rounds}), first-fit placement, the two-point throughput rule, penalty 1, and the ordering and
allocation named by the scheduler (\code{X+fair}, \code{X}, \code{X+tune}).

\subsection{Comparing}
For each configuration the 1{,}000 measured jobs are compared one by one: a job counts as exact if its JCT
differs by less than 1\,s. Means, 99th percentiles and the largest difference are also reported.

%======================================================================
\section{Results}
%======================================================================

\begin{longtable}{@{}llrrr@{}}
\toprule
Configuration & Synergy scheduler & Exact jobs & Mean JCT ours (h) & Synergy (h) \\
\midrule
\endhead
\multicolumn{5}{@{}l}{\emph{Step 0 (\code{presets/synergy-replica}): 9 jobs/hr, split (20, 70, 10), single-GPU}} \\
reference run & FIFO+tune & 1000/1000 & 21.8766 & 21.8766 \\
fixed-share baseline & FIFO+fair & 1000/1000 & 81.0886 & 81.0886 \\
\multicolumn{5}{@{}l}{\emph{Load sweep (paper Fig.~9), single-GPU}} \\
3 jobs/hr & FIFO+tune & 1000/1000 & 15.2479 & 15.2479 \\
3 jobs/hr & FIFO+fair & 1000/1000 & 17.3024 & 17.3024 \\
6 jobs/hr & FIFO+tune & 1000/1000 & 15.2480 & 15.2480 \\
6 jobs/hr & FIFO+fair & 1000/1000 & 17.3025 & 17.3025 \\
8 jobs/hr & FIFO+tune & 1000/1000 & 15.2527 & 15.2527 \\
8 jobs/hr & FIFO+fair & 1000/1000 & 27.4251 & 27.4251 \\
\multicolumn{5}{@{}l}{\emph{Allocation and orderings, 9 jobs/hr, single-GPU}} \\
Greedy & FIFO & 1000/1000 & 24.9841 & 24.9841 \\
SRTF & SRTF+tune & 1000/1000 & 16.6031 & 16.6031 \\
SRTF & SRTF+fair & 1000/1000 & 38.0276 & 38.0276 \\
LAS & LAS+tune & 1000/1000 & 33.1457 & 33.1457 \\
LAS & LAS+fair & 1000/1000 & 85.7921 & 85.7921 \\
FTF & FTF+tune & 1000/1000 & 366.4746 & 366.4746 \\
FTF & FTF+fair & 1000/1000 & 613.5850 & 613.5850 \\
\multicolumn{5}{@{}l}{\emph{CPU:GPU ratio (paper Fig.~12), 9 jobs/hr, profiles \code{models-4/5/6}}} \\
32 cores & FIFO+tune & 1000/1000 & 32.0849 & 32.0849 \\
32 cores & FIFO+fair & 1000/1000 & 81.0098 & 81.0098 \\
40 cores & FIFO+tune & 1000/1000 & 42.0332 & 42.0332 \\
40 cores & FIFO+fair & 1000/1000 & 81.0164 & 81.0164 \\
48 cores & FIFO+tune & 1000/1000 & 54.1124 & 54.1124 \\
48 cores & FIFO+fair & 1000/1000 & 81.0097 & 81.0097 \\
\multicolumn{5}{@{}l}{\emph{Workload split (paper Fig.~11), 9 jobs/hr}} \\
(33, 33, 33) & FIFO+tune & 1000/1000 & 31.0394 & 31.0394 \\
(33, 33, 33) & FIFO+fair & 1000/1000 & 81.0779 & 81.0779 \\
(33, 33, 33) & FIFO & 1000/1000 & 45.5326 & 45.5326 \\
(50, 0, 50) & FIFO+tune & 1000/1000 & 74.2724 & 74.2724 \\
(50, 0, 50) & FIFO+fair & 1000/1000 & 81.1122 & 81.1122 \\
(50, 0, 50) & FIFO & 1000/1000 & 158.8366 & 158.8366 \\
\multicolumn{5}{@{}l}{\emph{Multi-GPU jobs up to 16 GPUs (paper Figs.~7, 8), SOFT locality}} \\
4 jobs/hr & LAS+tune & 1000/1000 & 26.7152 & 26.7152 \\
4 jobs/hr & LAS+fair & 1000/1000 & 42.6008 & 42.6008 \\
5.5 jobs/hr & SRTF+tune & 1000/1000 & 44.1313 & 44.1313 \\
5.5 jobs/hr & SRTF+fair & 1000/1000 & 59.7479 & 59.7479 \\
4 jobs/hr & FIFO+tune & 1000/1000 & 37.0798 & 37.0798 \\
\bottomrule
\end{longtable}

In every row the largest difference of any job is 0.0\,s. The trends the paper reports are visible in
Synergy's own numbers and therefore in ours: Tune beats the fixed share at every load, by most at high load
(81.1\,h vs 21.9\,h at 9 jobs per hour, paper Fig.~9: 81\,h vs 22\,h); the gap shrinks as the CPU:GPU ratio
grows; Greedy collapses when most jobs are CPU-hungry (split (50, 0, 50): 158.8\,h, worse than the fixed share).

\paragraph{A note on the 32--48-core rows.} Tune's mean JCT rises from 21.9\,h (24 cores) to 54.1\,h (48 cores)
while the fixed share stays at 81\,h. This is not more CPUs making things worse: Synergy's \code{models-4/5/6}
profiles state each model's speedup relative to that cluster's own fixed share, so the same workload is a
different workload on each cluster (Phase~1c report). We reproduce Synergy here; interpreting Fig.~12 needs the
paper's own setup.

%======================================================================
\section{Bugs found and fixed}
%======================================================================

The first comparisons were not all exact. Each mismatch was traced to its cause; all three were in our module.

\subsection{An arrival exactly on a round boundary waited a round}
At 6 jobs per hour one job of 1{,}000 finished 300\,s late. It arrived at 2{,}447{,}700\,s, exactly on a round
boundary. Synergy processes events of equal time in the order arrival, job end, lease end, schedule
(\code{events/cluster\_event.py}), so the job joins that round. In CloudSim our round event had been queued before
the arrival event, so the round ran first. Synergy's arrival times are whole seconds, so about one job in 300
lands on a boundary. \textbf{Fix:} a round first admits every job whose arrival time has come; the job's own arrival
event is then ignored. A test reproduces the exact event order (it fails without the fix).

\subsection{FTF's starvation boost was missing}
FTF with Tune matched only 650 of 1{,}000 jobs (mean error 16\%). Synergy's FTF scheduler, like SRTF and LAS, gives a
job not run for 1000 hours priority from the next round on; in Phase~3 we had read its code with a truncated
search and concluded it did not, and wrote so in the code and the Phase~3 report. The run lasted 1{,}564
simulated hours, so the boost mattered. \textbf{Fix:} FTF uses the same starvation guard; the Phase~3 report is
corrected; a test covers it.

\subsection{LAS with Tune places the largest jobs first}
Multi-GPU LAS with Tune matched 887 of 1{,}000 jobs (mean error 1.0\%, inside the 10\% rule but not exact).
Synergy logs no allocations, so we added an opt-in plan trace to our broker (\code{-Dcosched.planTrace=<file>})
and the same dump to a copy of Synergy, and compared all 26{,}062 rounds. The first difference was at round 104:
Synergy placed an 8-GPU job before three 1-GPU jobs with lower LAS keys. Synergy's LAS scheduler, and no other,
chooses the jobs that fit the free GPU count in LAS order and then, when Tune is on, places them \emph{largest GPU
demand first} (\code{schedulers/las\_synergy\_new.py}). The paper does not describe this, and with single-GPU jobs
it changes nothing, which is why single-GPU LAS+tune matched. \textbf{Fix:} an explicit option on the ordering,
\code{largest\_first}, off by default and set for Synergy's LAS+tune. After the fix all 26{,}062 rounds are identical,
plan for plan.

%======================================================================
\section{Issues outside our module}
%======================================================================

\paragraph{Synergy's multi-GPU job generation is not reproducible.} Two runs of the same multi-GPU configuration
generated different arrivals despite the fixed seed, so Synergy's published multi-GPU numbers cannot be
regenerated exactly by rerunning it. Its \emph{scheduling} is deterministic: replaying a recorded trace with
\code{--replay\_trace} gives identical results. Our method is unaffected because every replay uses the trace
recorded by the same run it is compared with.

\paragraph{\code{tools/compare\_runs.py} compares means over different job sets.} It averages our file over all
jobs (6{,}925) but the reference over its 1{,}000, so an exact replay reports a 3.9\% error. Its job-by-job figures
are correct. Recorded as an open item for the tool's author.

%======================================================================
\section{Regression tests}
%======================================================================

\begin{itemize}
  \item \code{SynergyValidationTest} (every build, about 30\,s): the reference run and the fixed-share baseline from
        \code{presets/synergy-replica} must match Synergy job by job.
  \item \code{SynergyValidationSuiteTest} (opt-in): the 30 stored configurations in
        \code{presets/synergy-validation} (\(-\)\code{Dcosched.validate=subset} or \code{=all}); all 30 pass. Each holds the
        recorded trace, Synergy's JCTs and the experiment file; \code{tools/synergy\_validation} rebuilds them.
  \item Unit tests for each fix (boundary arrival, FTF starvation).
\end{itemize}
Project: 265 tests pass (one skipped: the opt-in suite).

%======================================================================
\section{Limits and what is not done}
%======================================================================
\begin{itemize}
  \item \textbf{What an exact match proves.} We reproduce Synergy's model of a cluster exactly, including its
        quirks. It does not show that the model is right about real clusters; for that the paper's physical runs
        (within 5\% of its simulator, \S5.2) are the evidence.
  \item \textbf{DRF and Tetris} (Fig.~13) are not implemented yet.
  \item \textbf{The real Philly arrival run} (Fig.~6, 512 GPUs, 8{,}000 jobs) was not run.
  \item \textbf{Part B}, checking our own Poisson job generator statistically against Synergy's, was not needed for
        validating scheduling, because we replay Synergy's own jobs; it matters only if we generate new traces.
  \item Exactness relies on \code{synergy\_rounds} (whole iterations per round, finishes at round ends). Our own
        experiments can use continuous progress; Phase~3 measured that this moves mean JCT by well under 1\%.
\end{itemize}

%======================================================================
\section{Next step}
%======================================================================
With the gate passed, Phase~5 runs our own experiments: preemption ablation, the Alibaba PAI robustness check,
the 2026 relevance check, and DRF and Tetris; and the comparison runner and decision report.

\end{document}
